\section{符号说明}

为避免四个问题分别定义不同符号，本文按“公共物理量—基础决策变量—后续新增业务变量”的层级统一记号，
全部符号集中在表~\ref{tab:symbol} 中说明。表内阴影行是分类标题，同一符号在后续问题中含义不变。

% 用 longtable 而非 table：符号较多，若一页放不下需要能自然续页，
% 表头在续页自动重复（table 环境放不下时会溢出到页外）。
\begin{longtable}{>{\raggedright\arraybackslash}p{2.3cm} >{\raggedright\arraybackslash}p{8.8cm} >{\raggedright\arraybackslash}p{3.2cm}}
\caption{全文符号说明}\label{tab:symbol}\\
\toprule
符号 & 含义 & 单位 / 来源 \\
\midrule
\endfirsthead
\multicolumn{3}{l}{\small 续表~\ref{tab:symbol}：全文符号说明}\\
\toprule
符号 & 含义 & 单位 / 来源 \\
\midrule
\endhead
\bottomrule
\endlastfoot

\multicolumn{3}{l}{\textit{（一）下标与集合}} \\
\(d\) & 日期索引 & 滚动计算中区分不同日期 \\
\(t\) & 日内第 \(t\) 个 10 min 时段，\(t=1,\ldots,144\) & 一天按 10 min 离散 \\
\(\tau\) & 滚动视野中的未来日期 & 区分当前日与未来日 \\
\(s\) & 日内调整阶段，\(s\in\{0,6,12,18\}\) & 预报发布时刻 \\
\(\omega\) & SAA 情景编号，\(\omega=1,\ldots,K\) & 由历史预测误差抽样 \\
\(K\) & SAA 情景数，正式取 \(K=4\) & 随机情景样本量 \\
\(R\) & 滚动视野长度，正式取 \(R=7\) & 天 \\
\midrule

\multicolumn{3}{l}{\textit{（二）已知参数与输入数据}} \\
\(L_{d,t}\) & 小区负荷功率 & kW \\
\(PV_{d,t}\) & 光伏发电功率 & kW \\
\(\pi_t\) & 固定日内电价（问题一至三） & 元/kWh \\
\(\pi^{act}_{d,t}\) & 问题四已实现的真实电价 & 元/kWh \\
\(\hat\pi_{d,t}\) & 计划阶段的电价中心预测 & 元/kWh \\
\(\pi^{(\omega,s)}_{d,t}\) & 阶段 \(s\) 下第 \(\omega\) 个价格情景 & 元/kWh \\
\(E_{\min},E_{\max}\) & 储能运行电量下限与上限，1200 与 10800 & kWh \\
\(P_{\max}\) & 最大充放电功率，5000 & kW \\
\(\eta_c,\eta_d\) & 充电效率、放电效率，均为 0.90 & — \\
\(\Delta t\) & 单个调度时段长度，\(1/6\) & h \\
\midrule

\multicolumn{3}{l}{\textit{（三）由输入数据预构造的辅助量（定义见问题一 5.2 节）}} \\
\(G\) & 外网总购电功率，\(G=G^L+G^{ch}\) & kW \\
\midrule

\multicolumn{3}{l}{\textit{（四）四问共用的基础物理决策变量}} \\
\(G^L\) & 正常外网购电中直接供负荷的功率 & kW \\
\(G^{ch}\) & 正常外网购电中用于储能充电的功率 & kW \\
\(PV^{ch}\) & 富余光伏用于储能充电的功率 & kW \\
\(C,D\) & 储能充电功率、放电功率 & kW \\
\(E\) & 时段结束后的储电量 & kWh \\
\(V\) & 弃光功率 & kW \\
\(u\) & 储能充放电状态二元变量 & 0/1 \\
\midrule

\multicolumn{3}{l}{\textit{（五）问题二新增变量}} \\
\(G^{plan}_{d,t}\) & 每天 0:00 提交的全天正常计划购电功率 & kW \\
\(H_{d,t}\) & 紧急购电功率（按 5 倍电价结算） & kW \\
\(W_{d,t}\) & 已购未用的正常电功率 & kW \\
\multicolumn{3}{l}{\quad 计划与执行之间满足 \(G^L+G^{ch}+W=G^{plan}\)} \\
\midrule

\multicolumn{3}{l}{\textit{（六）问题三新增变量}} \\
\(P_{d,t}\) & 0:00 原始计划购电功率 & kW \\
\(A^{s}_{d,t}\) & 阶段 \(s\) 调整后的购电功率 & kW \\
\(B_{d,t}\) & 全天最终生效计划（由四阶段生效区间拼接） & kW \\
\(\Delta^{+}_{d,t}\) & 相对原计划的最终调增量（按 1.5 倍电价结算） & kW \\
\(\Delta^{-}_{d,t}\) & 相对原计划的最终调减量（按 0.5 倍电价结算） & kW \\
\multicolumn{3}{l}{\quad 调整统一相对原计划结算：\(B-P=\Delta^{+}-\Delta^{-}\)} \\
\midrule

\multicolumn{3}{l}{\textit{（七）问题四新增变量}} \\
\(e^{\pi}\) & 历史价格预测残差 & 用于生成价格情景 \\
\(U_{d,t}\) & 计划购电物理上界，\(U_{d,t}=\max_{\omega}\overline L^{(\omega)}_{d,t}+P_{\max}\) & kW \\
\(z_{d,t}\) & 调增与调减的互斥二元变量 & 0/1 \\
\multicolumn{3}{l}{\quad 负价格下需显式互斥：\(0\le\Delta^{+}\le Uz\)，\(0\le\Delta^{-}\le U(1-z)\)} \\

\end{longtable}

变量体系逐层继承而非每问重新定义：问题一引入基础物理变量，
问题二在其上加入计划与紧急购电相关变量，问题三再加入阶段化的计划调整变量，
问题四最后加入价格不确定性与套利防范变量。高编号问题的新增变量是对前序问题的
扩展，低编号问题的符号含义在后文中保持不变。
